% 3.2.2 硬件映射图；基准功率计量组与热节点分开。
\begin{figure}[htbp]
\centering
\resizebox{0.98\textwidth}{!}{%
\begin{tikzpicture}[font=\small,>=Latex,
 unit/.style={draw=blue!60!black,fill=blue!6,rounded corners=2pt,minimum height=.63cm,minimum width=1.43cm,align=center},
 freqbox/.style={draw=green!55!black,fill=green!8,rounded corners=2pt,minimum height=.72cm,align=center},
 power/.style={draw=orange!70!black,fill=orange!8,rounded corners=2pt,minimum height=.72cm,align=center},
 thermal/.style={draw=red!55!black,fill=red!6,rounded corners=2pt,minimum height=.65cm,minimum width=1.43cm,align=center}]
 \node[anchor=east,font=\bfseries] at (-.45,4.35) {执行单元};
 \node[anchor=east,font=\bfseries] at (-.45,2.98) {控制域};
 \node[anchor=east,font=\bfseries] at (-.45,1.57) {计量组};
 \node[anchor=east,font=\bfseries] at (-.45,.15) {热节点};
 \node[unit] (cu1) at (1.22,4.35) {CPU1};
 \node[unit] (cu2) at (3.05,4.35) {CPU2};
 \node[unit] (gu) at (6.35,4.35) {GPU0};
 \node[unit] (nu) at (9.90,4.35) {NPU0};
 \node[freqbox,minimum width=3.74cm] (cd) at (2.13,2.98) {CPU共享频率域};
 \node[freqbox,minimum width=2.3cm] (gd) at (6.35,2.98) {GPU域};
 \node[freqbox,minimum width=2.3cm] (nd) at (9.90,2.98) {NPU域};
 \node[power,minimum width=3.74cm] (cp) at (2.13,1.57) {CPU基准计费一次};
 \node[power,minimum width=2.3cm] (gp) at (6.35,1.57) {GPU基准};
 \node[power,minimum width=2.3cm] (np) at (9.90,1.57) {NPU基准};
 \node[thermal] (ct1) at (1.22,.15) {CPU1节点};
 \node[thermal] (ct2) at (3.05,.15) {CPU2节点};
 \node[thermal] (gt) at (6.35,.15) {GPU节点};
 \node[thermal] (nt) at (9.90,.15) {NPU节点};
 \foreach \a/\b in {cu1/cd,cu2/cd,gu/gd,nu/nd,cd/cp,gd/gp,nd/np,cp/ct1,cp/ct2,gp/gt,np/nt}
  {\draw[-Latex,gray!65,thick] (\a)--(\b);}
 \node[font=\scriptsize,align=center,red!55!black] at (2.14,-.58) {本地温度分开计算\\CPU基准按输入分配注热};
 \node[font=\scriptsize] at (8.22,-.38) {不同域可独立选择工作点};
\end{tikzpicture}}
\caption{执行单元、频率控制域、功率计量组和热节点的不同粒度}
\label{fig:ch3-hardware-layers}
\end{figure}
